\subsection{A binary trivial leaf at a five-element center}

\begin{theorem}[Proved: a five-element center with a binary trivial leaf]
\label{CP2:A-five-center-binary-trivial-leaf-nineteen}
Let $C$ be a three-object source star with at most nineteen arrows.
Suppose its center has order five, one leaf is trivial with a
two-element stalk, and the other leaf has order at most eleven.
Then $BC\not\simeq\CP^2$. The same conclusion holds for the
opposite category.
\end{theorem}
\begin{proof}
Write $M$ for the center, $S$ for the binary right $M$-set at
the trivial leaf, and $N,T$ for the other leaf and stalk.
Assume for contradiction that $BC\simeq\CP^2$.

If $S$ has an $M$-fixed point, collapse $S$ to that point and
leave $T$ unchanged. This is an idempotent natural endomorphism
of the connecting profunctor from the center to the disjoint
union of the leaves. Theorem~\ref{CP2:A-cut-retract} transfers
actual projective-plane type to its image. That image has one
fewer arrow, contradicting Theorem~\ref{CP2:A-global-lower-nineteen}.
We may assume that $S$ has no fixed point.

Let $K$ be the minimal ideal of $M$, with minimal group $G$.
If $M\ne P_5$ and $G$ is trivial, the actual small-minimal-ideal
comparison gives $BM\simeq *$, by
Corollary~\ref{CP2:A-five-types}. The full star nerve is the
union of its two edge nerves along $BM$, hence is homotopy
equivalent to the wedge of those two edge spaces. A wedge of
connected CW spaces having type $\CP^2$ has one wedge summand
of that same type: the fundamental groups are free factors of
the trivial group, integral homology splits, and the primitive
second square belongs to exactly one summand. Both edge
categories have at most seventeen arrows, since deleting the
other nonempty stalk and its leaf removes at least two arrows.
This again contradicts the established eighteen-arrow exclusion.

If $K$ has one row, the trivial left $\mathbb QM$-module is
projective, by Lemma~\ref{CP2:A-projectivity-test}. The actual
right-orbit quotient replacing the center by the trivial monoid
is a rational homology and cohomology-ring equivalence, as in
Lemma~\ref{CP2:A-high-orbits}. Its star is a wedge of the ordinary
act cones on $S/M$ and $T/M$, at leaf corners $1$ and $N$.
Their rational homology is bounded, since it is a direct
summand of that of the quotient star. The bounded-eleven act
theorem makes every second product in both summands zero,
and hence in the star, a contradiction.

The remaining five-element possibilities are $P_5$ and
\[
 L=(\{0,1\}\times C_2)^1,
 \qquad (i,g)(j,h)=(i,gh).
\]
Indeed, if $G$ is nontrivial and there are at least two rows,
the ideal already has at least four elements. The order-five
budget forces exactly two rows, one column, $G=C_2$, and only
the global identity outside $K$. If $G$ is trivial, the only
case not already treated is $P_5$.

For either remaining monoid, every fixed-point-free binary
right act is a representable minimal right ideal. For $P_5$,
the band actors are idempotent transformations of two points.
If they are all permutations, they are identities, which gives
fixed points. Otherwise the image of the minimal ideal consists
of constant maps; absence of a fixed point forces both constants
to occur. The band equation makes their values depend only on
the column, and the two column values are distinct. Thus
$S\cong eM$ for any minimal idempotent $e$, by choosing $s$ with
$se=s$ and sending $em$ to $sm$.

For $L$, the core $SK$ is a union of one- or two-element
$C_2$ orbits. A singleton core orbit is fixed by every element
of $K$ and by the outside identity, so it is excluded.
Therefore $SK=S$ is one regular $C_2$ orbit. A minimal right
ideal $eM$ is the regular $C_2$ right act, and the same map
$em\mapsto sm$ is again a right-$M$-equivariant bijection.

The action category of $eM$ is an actual retract of the regular
right action category of $M$, and the latter is contractible.
Hence $B(S,M)$ is contractible. The binary edge space is the
homotopy pushout
\[
 BM\ \longleftarrow\ B(S,M)\ \longrightarrow\ *,
\]
so its inclusion of $BM$ is a homotopy equivalence. Gluing the
other edge along that same corner shows that deleting the
binary stalk and its trivial leaf preserves the homotopy type
of the entire star. The remaining category has at most sixteen
arrows, contradicting the eighteen-arrow exclusion once more.
\end{proof}

\begin{corollary}[Proved: the five-center large-leaf frontier]
\label{CP2:A-star-five-one-eight-two-three}
The nineteen-arrow star profile $(5,1,8;2,3)$ cannot realize
$\CP^2$.
\end{corollary}
\begin{proof}
It satisfies the preceding theorem.
\end{proof}
